% Options for packages loaded elsewhere
\PassOptionsToPackage{unicode}{hyperref}
\PassOptionsToPackage{hyphens}{url}
\PassOptionsToPackage{dvipsnames,svgnames,x11names}{xcolor}
%
\documentclass[
  11pt,
]{article}
\usepackage{amsmath,amssymb}
\usepackage{iftex}
\ifPDFTeX
  \usepackage[T1]{fontenc}
  \usepackage[utf8]{inputenc}
  \usepackage{textcomp} % provide euro and other symbols
\else % if luatex or xetex
  \usepackage{unicode-math} % this also loads fontspec
  \defaultfontfeatures{Scale=MatchLowercase}
  \defaultfontfeatures[\rmfamily]{Ligatures=TeX,Scale=1}
\fi
\usepackage[]{lmodern}
\ifPDFTeX\else
  % xetex/luatex font selection
\fi
% Use upquote if available, for straight quotes in verbatim environments
\IfFileExists{upquote.sty}{\usepackage{upquote}}{}
\IfFileExists{microtype.sty}{% use microtype if available
  \usepackage[]{microtype}
  \UseMicrotypeSet[protrusion]{basicmath} % disable protrusion for tt fonts
}{}
\makeatletter
\@ifundefined{KOMAClassName}{% if non-KOMA class
  \IfFileExists{parskip.sty}{%
    \usepackage{parskip}
  }{% else
    \setlength{\parindent}{0pt}
    \setlength{\parskip}{6pt plus 2pt minus 1pt}}
}{% if KOMA class
  \KOMAoptions{parskip=half}}
\makeatother
\usepackage{xcolor}
\usepackage{longtable,booktabs,array}
\usepackage{calc} % for calculating minipage widths
% Correct order of tables after \paragraph or \subparagraph
\usepackage{etoolbox}
\makeatletter
\patchcmd\longtable{\par}{\if@noskipsec\mbox{}\fi\par}{}{}
\makeatother
% Allow footnotes in longtable head/foot
\IfFileExists{footnotehyper.sty}{\usepackage{footnotehyper}}{\usepackage{footnote}}
\makesavenoteenv{longtable}
\setlength{\emergencystretch}{3em} % prevent overfull lines
\providecommand{\tightlist}{%
  \setlength{\itemsep}{0pt}\setlength{\parskip}{0pt}}
\setcounter{secnumdepth}{-\maxdimen} % remove section numbering
\usepackage[a4paper,margin=2.2cm,headheight=14pt]{geometry}
\usepackage{fancyhdr}
\usepackage{titlesec}
\usepackage{xcolor}
\usepackage{microtype}
\definecolor{accent}{HTML}{1F4E79}
\titleformat{\section}{\large\bfseries\color{accent}}{}{0em}{}
\titleformat{\subsection}{\normalsize\bfseries\color{accent}}{}{0em}{}
\pagestyle{fancy}\fancyhf{}
\lhead{\small Architecture Decision Records}\rhead{\small Modular Learner Experience (P10-G37)}\cfoot{\small\thepage}
\renewcommand{\arraystretch}{1.2}
\sloppy
\usepackage{newunicodechar}
\newunicodechar{≤}{\ensuremath{\le}}
\newunicodechar{≥}{\ensuremath{\ge}}
\newunicodechar{→}{\ensuremath{\rightarrow}}
\newunicodechar{·}{\textperiodcentered}
\newunicodechar{×}{\ensuremath{\times}}
\ifLuaTeX
  \usepackage{selnolig}  % disable illegal ligatures
\fi
\IfFileExists{bookmark.sty}{\usepackage{bookmark}}{\usepackage{hyperref}}
\IfFileExists{xurl.sty}{\usepackage{xurl}}{} % add URL line breaks if available
\urlstyle{same}
\hypersetup{
  pdftitle={Architecture Decision Records --- P10-G37},
  colorlinks=true,
  linkcolor={accent},
  filecolor={Maroon},
  citecolor={Blue},
  urlcolor={accent},
  pdfcreator={LaTeX via pandoc}}

\title{Architecture Decision Records --- P10-G37}
\author{}
\date{Written 21 November 2026}

\begin{document}
\maketitle

{
\hypersetup{linkcolor=accent}
\setcounter{tocdepth}{1}
\tableofcontents
}
Project 10 · Group 37 · Shubham Yadav · Repository:
https://github.com/shhubbhaam/agenticUI

\textbf{About the dates.} This log was written on \textbf{21 Nov 2026}.
Each decision date below comes from the commit or document that first
shows the decision; none is backdated. Records are never deleted: a
replaced decision is marked \emph{Superseded} and points to its
successor.

\begin{longtable}[]{@{}
  >{\raggedright\arraybackslash}p{(\columnwidth - 4\tabcolsep) * \real{0.3333}}
  >{\raggedright\arraybackslash}p{(\columnwidth - 4\tabcolsep) * \real{0.3333}}
  >{\raggedright\arraybackslash}p{(\columnwidth - 4\tabcolsep) * \real{0.3333}}@{}}
\toprule\noalign{}
\begin{minipage}[b]{\linewidth}\raggedright
ID
\end{minipage} & \begin{minipage}[b]{\linewidth}\raggedright
Title
\end{minipage} & \begin{minipage}[b]{\linewidth}\raggedright
Status
\end{minipage} \\
\midrule\noalign{}
\endhead
\bottomrule\noalign{}
\endlastfoot
ADR-001 & LearnHouse (team fork) is the LMS and source of truth &
Accepted \\
ADR-002 & A thin BFF holds all secrets, sessions and policy &
Accepted \\
ADR-003 & Agent selects from a trusted catalogue; one shared Zod
contract & Accepted \\
ADR-004 & Methodology record (9 Oct reaffirmation) & Accepted \\
ADR-005 & SQLite outbox + durable BFF inbox; conflict rules by data type
& Accepted (supersedes in-memory idempotency map) \\
ADR-006 & Dependency-free BFF on \texttt{node:http} and
\texttt{node:sqlite} & Accepted \\
ADR-007 & Checkpoints as a BFF extension with server grading &
Accepted \\
ADR-008 & AI coach over SSE (AG-UI-style subset) with bound approvals &
Accepted \\
ADR-009 & M3 web modules built inside the LearnHouse fork & Accepted
(approved scope change) \\
ADR-010 & Compile-time mobile feature switches & Accepted \\
ADR-011 & Text size follows Android's system setting & Accepted \\
ADR-012 & Video: YouTube player plus hosted uploads & Accepted \\
ADR-013 & Gemini only behind the BFF, enum-restricted structured output
& Accepted \\
\end{longtable}

\begin{center}\rule{0.5\linewidth}{0.5pt}\end{center}

\hypertarget{adr-001-learnhouse-team-fork-is-the-lms-and-source-of-truth}{%
\section{ADR-001 --- LearnHouse (team fork) is the LMS and source of
truth}\label{adr-001-learnhouse-team-fork-is-the-lms-and-source-of-truth}}

\textbf{Date, status, scope.} Late Sept 2026 (W1 review; adapter code
from 19 Sept, commit \texttt{0f62868}). Accepted. All milestones; M5
read-back.

\textbf{Context and drivers.} The project needs courses, learners,
enrolment and progress, but the course reviewers said not to build a new
LMS. The team already runs a LearnHouse fork at
\texttt{learn.praxisacademy.ai}.

\textbf{Options.} (a) Build a small LMS/mock backend. (b) Use LearnHouse
through its API (chosen). (c) Use LearnHouse's own web UI only.

\textbf{Decision and rationale.} (b): LearnHouse owns courses,
activities, enrolment, completion and grades; the app and BFF read and
write through its API and always read back. Judgement, not measurement:
this avoids a second course database and matches reviewer advice.

\textbf{Consequences.} Real routes must be pinned and verified (the
public OpenAPI is disabled, so the deployed commit is still unpinned).
Missing features (checkpoints, event inbox, agent sessions) need small
extension records (ADR-005, 007, 008).

\textbf{Evidence.} \texttt{bff/src/learnhouse.ts},
\texttt{bff/test/learnhouse.test.ts}; live run
\texttt{docs/evidence/live-run-2026-10-24.md}.

\textbf{Revisit trigger.} LearnHouse lacks a required write
(e.g.~per-checkpoint progress) or changes routes between releases.

\begin{center}\rule{0.5\linewidth}{0.5pt}\end{center}

\hypertarget{adr-002-a-thin-bff-holds-all-secrets-sessions-and-policy}{%
\section{ADR-002 --- A thin BFF holds all secrets, sessions and
policy}\label{adr-002-a-thin-bff-holds-all-secrets-sessions-and-policy}}

\textbf{Date, status, scope.} 11 Oct 2026
(\texttt{docs/W3-plan-and-evidence.md}). Accepted. M1, M5, M7, security.

\textbf{Context.} The W2 review found a Gemini key literal in the app
and the LearnHouse organisation token being used from the mobile bundle,
with learner ids taken from URLs.

\textbf{Options.} (a) Keep calls in the app with environment variables.
(b) A thin backend-for-frontend (chosen). (c) LearnHouse learner tokens
only, no backend.

\textbf{Decision.} (b): the BFF checks the login against LearnHouse,
issues an opaque expiring session, derives the learner from it on every
\texttt{/api/me/*} route, checks the course allow-list and enrolment,
and is the only holder of \texttt{LH\_API\_TOKEN} and
\texttt{GEMINI\_API\_KEY}. (c) was rejected because the progress write
route used needs the org token.

\textbf{Consequences.} One more process to run; sessions are in memory,
so a BFF restart signs learners out. The Gemini key committed before W3
was revoked on \texttt{{[}date\ —\ fill\ in{]}}; it remains in git
history.

\textbf{Evidence.} \texttt{bff/src/app.ts},
\texttt{bff/src/sessions.ts}, \texttt{bff/test/bff.test.ts} (21 tests),
\texttt{docs/evidence/secret-scan.txt}.

\textbf{Revisit trigger.} A second client needs sessions shared across
BFF instances, or LearnHouse exposes learner-scoped progress writes.

\begin{center}\rule{0.5\linewidth}{0.5pt}\end{center}

\hypertarget{adr-003-agent-selects-from-a-trusted-catalogue-one-shared-zod-contract}{%
\section{ADR-003 --- Agent selects from a trusted catalogue; one shared
Zod
contract}\label{adr-003-agent-selects-from-a-trusted-catalogue-one-shared-zod-contract}}

\textbf{Date, status, scope.} Proposal (W0), strict versioned contract
11 Oct 2026, server-side use of the same schema 14 Nov 2026
(\texttt{9fd4608}). Accepted. M1, M7.

\textbf{Context.} Agent output must never execute; the W3 review found
the BFF and app validating differently.

\textbf{Options.} (a) Rule-based layout only. (b) Agent-generated UI
code. (c) Agent chooses catalogue components and props; both sides
validate with one schema (chosen).

\textbf{Decision.} (c): four components, three typed actions, fixed
slots, versioned envelope with learner/course/expiry, UTF-8 byte cap,
per-slot fallback, error boundary per slot; the BFF runs the app's own
\texttt{parseHomeLayout} and drops slots the app would reject. (a) is
kept as the fallback when the model fails.

\textbf{Consequences.} Adding a component needs code, schema, registry
entry and tests. A valid layout is not proven to be a good one; no
relevance evaluation was run.

\textbf{Evidence.} \texttt{src/contracts/homeLayout.ts}, contract suite,
\texttt{bff/test/agent.test.ts}, \texttt{SlotBoundary.test.tsx}.

\textbf{Revisit trigger.} The catalogue grows beyond what can be
described in one prompt, or a reviewer requires official A2UI
conformance.

\begin{center}\rule{0.5\linewidth}{0.5pt}\end{center}

\hypertarget{adr-004-methodology-record-9-october-reaffirmation}{%
\section{ADR-004 --- Methodology record (9 October
reaffirmation)}\label{adr-004-methodology-record-9-october-reaffirmation}}

\textbf{Date, status, scope.} Methodology document dated \textbf{11 Oct
2026} --- two days after the 9 Oct 2026, 23:59 IST deadline. This record
written 21 Nov 2026. Accepted. All milestones.

\textbf{Statement.} The approach was \textbf{reaffirmed}, not changed,
in the W3 plan: hybrid catalogue-based agent UI over LearnHouse through
a thin BFF, with M1--M9 mapped to binary test groups T1.1--T9.2 and
thresholds (cold start ≤ 3 s, web bundle ≤ 500 KB gzip,
reconnect-to-synced ≤ 5 s, resume drift ≤ 1 s). The 13 scenarios were
mapped later (\texttt{docs/W5-status-and-evidence.md} §3).

\textbf{Later changes (not part of the 9 Oct record).}

\begin{longtable}[]{@{}
  >{\raggedright\arraybackslash}p{(\columnwidth - 8\tabcolsep) * \real{0.2000}}
  >{\raggedright\arraybackslash}p{(\columnwidth - 8\tabcolsep) * \real{0.2000}}
  >{\raggedright\arraybackslash}p{(\columnwidth - 8\tabcolsep) * \real{0.2000}}
  >{\raggedright\arraybackslash}p{(\columnwidth - 8\tabcolsep) * \real{0.2000}}
  >{\raggedright\arraybackslash}p{(\columnwidth - 8\tabcolsep) * \real{0.2000}}@{}}
\toprule\noalign{}
\begin{minipage}[b]{\linewidth}\raggedright
Old
\end{minipage} & \begin{minipage}[b]{\linewidth}\raggedright
Evidence prompting change
\end{minipage} & \begin{minipage}[b]{\linewidth}\raggedright
Alternatives
\end{minipage} & \begin{minipage}[b]{\linewidth}\raggedright
Selected
\end{minipage} & \begin{minipage}[b]{\linewidth}\raggedright
Impact
\end{minipage} \\
\midrule\noalign{}
\endhead
\bottomrule\noalign{}
\endlastfoot
M3 as Module-Federation shell + 3 remotes & Single member; web work
already happening in the LearnHouse fork & Separate React shell;
headless examples & Web features in the LearnHouse fork (ADR-009) &
T3.2/T3.3 (independent deploy, failure isolation) not demonstrated \\
Checkpoints in the app & LearnHouse has no checkpoints; answers on the
device are insecure & App-only; LearnHouse plugin & BFF extension
(ADR-007) & T4.1 tested on the BFF \\
First durable offline event in W3 & Not delivered in W3 & --- &
Delivered 14 Nov (ADR-005) & M5 later than planned \\
\end{longtable}

\textbf{Approval.} The M3 change was agreed with the reviewer in
discussion on \texttt{{[}date\ —\ fill\ in{]}}; this agreement is not
recorded elsewhere in the repository.

\begin{center}\rule{0.5\linewidth}{0.5pt}\end{center}

\hypertarget{adr-005-sqlite-outbox-durable-bff-inbox-conflict-rules-by-data-type}{%
\section{ADR-005 --- SQLite outbox + durable BFF inbox; conflict rules
by data
type}\label{adr-005-sqlite-outbox-durable-bff-inbox-conflict-rules-by-data-type}}

\textbf{Date, status, scope.} 14 Nov 2026 (\texttt{ad16a12}). Accepted;
supersedes the W3 in-memory \texttt{Idempotency-Key} map. M5.

\textbf{Context.} The W3 review showed that concurrent duplicate
requests could each write upstream, that the map was lost on restart,
and that ``received'' was shown as ``saved''.

\textbf{Options.} (a) Keep online-only writes with retries. (b)
AsyncStorage queue. (c) SQLite outbox on the device and a SQLite inbox
on the BFF (chosen).

\textbf{Decision.} Each event gets an id once and keeps it through
retries, plus device id, per-device sequence and schema version. The BFF
stores id + payload hash; same id/same payload returns the stored
result, same id/new payload is rejected. Results are \texttt{applied}
(read-back confirmed), \texttt{received} (retry later) or
\texttt{rejected} (kept visible). Completion is monotonic; quiz hand-ins
are append-only attempts with stage markers so a resend never submits
twice.

\textbf{Consequences.} Preference conflicts (S9) are not covered. The
device must keep rejected rows until the learner sees them.

\textbf{Evidence.} \texttt{src/data/local/outbox.ts},
\texttt{src/data/sync/syncEngine.ts}, \texttt{bff/src/inbox.ts},
\texttt{bff/test/events.test.ts} (13); device run 14 Nov 2026.

\textbf{Revisit trigger.} A second client type writes the same records
offline, or a multi-instance BFF is needed.

\begin{center}\rule{0.5\linewidth}{0.5pt}\end{center}

\hypertarget{adr-006-dependency-free-bff-on-nodehttp-and-nodesqlite}{%
\section{\texorpdfstring{ADR-006 --- Dependency-free BFF on
\texttt{node:http} and
\texttt{node:sqlite}}{ADR-006 --- Dependency-free BFF on node:http and node:sqlite}}\label{adr-006-dependency-free-bff-on-nodehttp-and-nodesqlite}}

\textbf{Date, status, scope.} 11 Oct 2026. Accepted. BFF.

\textbf{Context.} The development sandbox could not install npm
packages; Express/axios/dotenv were in use.

\textbf{Options.} (a) Keep Express and install locally only. (b) Node
built-ins: \texttt{node:http}, native \texttt{-\/-env-file},
\texttt{node:sqlite}, \texttt{node:test} (chosen).

\textbf{Decision.} (b), with a small hand-written router.
\textbf{Consequences.} No middleware ecosystem; edge cases (malformed
URI) must be handled by hand. \textbf{Evidence.}
\texttt{bff/src/app.ts}; 121/121 tests on Node 22.23.5 (21 Nov 2026).

\textbf{Revisit trigger.} Routing or validation code starts growing
faster than features.

\begin{center}\rule{0.5\linewidth}{0.5pt}\end{center}

\hypertarget{adr-007-checkpoints-as-a-bff-extension-with-server-grading}{%
\section{ADR-007 --- Checkpoints as a BFF extension with server
grading}\label{adr-007-checkpoints-as-a-bff-extension-with-server-grading}}

\textbf{Date, status, scope.} First attempt in the app 6 Nov
(\texttt{3ca95ca}); moved to the BFF 6 Nov (\texttt{a780226}) and
finished 14 Nov (\texttt{37f9cc3}). Accepted. M4.

\textbf{Options.} (a) Checkpoints and answers in the app. (b) A
LearnHouse plugin. (c) BFF file + store (chosen).

\textbf{Decision.} \texttt{bff/checkpoints.json} (versioned, per
activity or \texttt{*}); the app receives questions without answers; the
BFF grades, appends attempts, enforces order and refuses LearnHouse
completion until all are passed.

\textbf{Consequences.} Checkpoint content is synthetic and shared by
every video; authors cannot edit it in LearnHouse.

\textbf{Evidence.} \texttt{bff/test/checkpoints.test.ts} (10),
\texttt{src/screens/lesson/lesson.test.tsx}.

\textbf{Revisit trigger.} Authors need per-lesson checkpoints managed in
the LMS.

\begin{center}\rule{0.5\linewidth}{0.5pt}\end{center}

\hypertarget{adr-008-ai-coach-over-sse-ag-ui-style-subset-with-bound-approvals}{%
\section{ADR-008 --- AI coach over SSE (AG-UI-style subset) with bound
approvals}\label{adr-008-ai-coach-over-sse-ag-ui-style-subset-with-bound-approvals}}

\textbf{Date, status, scope.} 14 Nov 2026 (\texttt{3c032e8}); redirect
after approval 15 Nov (\texttt{1571c7f} + uncommitted
\texttt{CoachScreen.tsx}). Accepted. M6, S4, S5, S10, S13.

\textbf{Options.} (a) WebSocket. (b) Official AG-UI SDK. (c) SSE with a
documented AG-UI-style event subset (chosen).

\textbf{Decision.} SSE over \texttt{expo/fetch}; event ids = sequence;
\texttt{Last-Event-ID} resume; 5-min / 500-event replay window with a
snapshot outside it; 15 s heartbeat; client back-off and 45 s stall
watchdog. Gemini chooses only among server-offered options. A proposal
has a SHA-256 hash, a 10-min expiry, one decision, and is bound to the
learner and login session; the target is re-checked when the decision
runs. The coach is isolated from the Home layout.

\textbf{Consequences.} Not full AG-UI conformance; sessions are lost on
BFF restart; web streaming not built.

\textbf{Evidence.} \texttt{bff/src/coach.ts},
\texttt{bff/test/coach.test.ts} (13), \texttt{src/core/coach/*.test.ts}.

\textbf{Revisit trigger.} A second client needs the stream, or approvals
must survive a BFF restart.

\begin{center}\rule{0.5\linewidth}{0.5pt}\end{center}

\hypertarget{adr-009-m3-web-modules-built-inside-the-learnhouse-fork}{%
\section{ADR-009 --- M3 web modules built inside the LearnHouse
fork}\label{adr-009-m3-web-modules-built-inside-the-learnhouse-fork}}

\textbf{Date, status, scope.} Work from 12 Oct 2026 (fork
\texttt{4adc2055}) to 21 Nov 2026 (\texttt{6a301103}); agreement with
reviewer \texttt{{[}date\ —\ fill\ in{]}}. Accepted as an approved scope
change. M3.

\textbf{Options.} (a) Module-Federation shell with three remote modules
over the BFF. (b) Features in the LearnHouse web app (chosen).

\textbf{Decision.} Three modules --- community image upload, Help \&
Support tickets, weighted course grading --- reuse LearnHouse auth,
routing, RBAC, design system and i18n.

\textbf{Consequences.} Not independently built or deployed; no
module-failure isolation test; they do not consume the BFF contracts (M1
web consumer missing). Their 131 API and 37 web tests were not run for
the final report. LearnHouse is AGPL-3.0, so the deployed fork's source
must be offered to its users.

\textbf{Evidence.} https://github.com/shhubbhaam/learnhouse commits
\texttt{4adc2055}, \texttt{d39e16c7}, \texttt{1108064b},
\texttt{04829b8f}, \texttt{6a301103}.

\textbf{Revisit trigger.} Evaluators require T3.2/T3.3 evidence as
originally specified.

\begin{center}\rule{0.5\linewidth}{0.5pt}\end{center}

\hypertarget{adr-010-compile-time-mobile-feature-switches}{%
\section{ADR-010 --- Compile-time mobile feature
switches}\label{adr-010-compile-time-mobile-feature-switches}}

\textbf{Date, status, scope.} 14 Nov 2026 (\texttt{9fd4608}). Accepted.
M2.

\textbf{Options.} (a) Runtime-loaded mobile modules (excluded by W0
advice). (b) Separate apps. (c) Compile-time switches per module
(chosen).

\textbf{Decision.} \texttt{src/config/features.ts} turns
\texttt{player}, \texttt{assessment}, \texttt{agentUi} off
independently; their routes open a ``turned off'' screen; Home falls
back to progress cards. \textbf{Evidence.}
\texttt{src/navigation/activityRoute.test.ts}.

\textbf{Revisit trigger.} A module must be updated without an app
release.

\begin{center}\rule{0.5\linewidth}{0.5pt}\end{center}

\hypertarget{adr-011-text-size-follows-androids-system-setting}{%
\section{ADR-011 --- Text size follows Android's system
setting}\label{adr-011-text-size-follows-androids-system-setting}}

\textbf{Date, status, scope.} 15 Nov 2026 (\texttt{114f77e}). Accepted.
M8.

\textbf{Options.} (a) In-app ``Larger text'' multiplier. (b) Follow the
OS font scale (chosen), keep the in-app switch untouched for now.

\textbf{Decision.} Font scaling is never disabled; fixed-height
containers became \texttt{minHeight}; only text in fixed grids or bars
(streak calendar, tab labels, player times, study-time tiles) is capped
at 1.3--1.5×.

\textbf{Consequences.} The existing in-app ``Larger text'' switch has no
effect yet. \textbf{Evidence.} \texttt{src/ui/fontScaling.test.tsx};
200\% run observed once, not recorded.

\textbf{Revisit trigger.} TalkBack or 200\% recordings show clipped
content.

\begin{center}\rule{0.5\linewidth}{0.5pt}\end{center}

\hypertarget{adr-012-video-youtube-player-plus-hosted-uploads}{%
\section{ADR-012 --- Video: YouTube player plus hosted
uploads}\label{adr-012-video-youtube-player-plus-hosted-uploads}}

\textbf{Date, status, scope.} 31 Oct 2026 (\texttt{730f7cf}). Accepted.
M4.

\textbf{Decision.} Course videos play through the YouTube player
(supervisor's direction); at least one video uploaded to LearnHouse
plays through \texttt{expo-video} via the BFF file route.
\textbf{Consequences.} Captions rely on YouTube's own; uploaded videos
have none.

\begin{center}\rule{0.5\linewidth}{0.5pt}\end{center}

\hypertarget{adr-013-gemini-only-behind-the-bff-enum-restricted-structured-output}{%
\section{ADR-013 --- Gemini only behind the BFF, enum-restricted
structured
output}\label{adr-013-gemini-only-behind-the-bff-enum-restricted-structured-output}}

\textbf{Date, status, scope.} 11 Oct 2026; coach model 14 Nov 2026.
Accepted. M6, M7, S12.

\textbf{Decision.} Calls go only from the BFF with a response schema
whose ids are an enum of server-authorised activities or options; output
is validated again; template/deterministic fallback on any failure; only
the chosen id is logged, never model text; the model gets no name,
e-mail or raw answers.

\textbf{Evidence.} \texttt{bff/src/agent.ts},
\texttt{bff/src/coachModel.ts}, \texttt{bff/test/coach.test.ts} (model
output outside the options; no model text in logs).

\textbf{Revisit trigger.} Model latency (9.4 s observed once on 24 Oct)
becomes the main UX problem, or another provider is required.

\end{document}
